\input{preamble}

\begin{document}

\title{Linguistics}
\maketitle

\phantomsection
\label{section-phantom}

\tableofcontents

\section{Ancient Greek orthography}

\subsection{Breathings}

In polytonic Ancient Greek, a word beginning with a vowel normally carries a
\emph{breathing}.  The \emph{smooth breathing} (Greek \emph{psili}) marks the
absence of an initial consonantal sound, whereas the \emph{rough breathing}
(Greek \emph{daseia}) marks an initial /h/.  Thus a form transliterated
\emph{e-} with smooth breathing begins directly with the vowel, while a form
transliterated \emph{he-} with rough breathing begins with a light expiration
/h/.  The rough breathing is therefore an \emph{expiration}, not an
inspiration.

For example, the Homeric verb \emph{helisso}, ``to turn, twist, or curve'',
begins with rough breathing and is transliterated with an initial \emph{h}.
This is the verb behind the Homeric epithet \emph{amphielissai}, used of ships
and traditionally understood as ``curved on both sides'' or ``double-curved''.

A breathing is not an accent.  On an initial vowel, a breathing and an accent
may occur simultaneously.

\subsection{Accents}

The three traditional accents are the acute, grave, and circumflex.  In the
classical pitch-accent system they concern pitch rather than nasalization or
mere stress.

The acute marks a high pitch.  The circumflex can occur only on a long
syllable and represents a high--low contour within that syllable.  The grave
replaces a final acute in certain contexts when another word follows.  The
circumflex is sometimes printed in a shape resembling a tilde or a caret, but
it has nothing to do with the Portuguese tilde: it does not nasalize the
vowel.  A long vowel bearing a circumflex is long independently of the
circumflex; the accent indicates its pitch contour.

\section{Ancient Greek verbs}

\subsection{Regular present active endings}

A useful first model for the present indicative active is
\begin{center}
\begin{tabular}{c|cc}
 & Singular & Plural \\
\hline
1st person & \texttt{-o} & \texttt{-omen} \\
2nd person & \texttt{-eis} & \texttt{-ete} \\
3rd person & \texttt{-ei} & \texttt{-ousi(n)}
\end{tabular}
\end{center}
Thus, for example, \emph{luo} ``I release'' gives \emph{luo, lueis, luei,
luomen, luete, luousi(n)}.

The final \emph{-n} in forms such as \emph{-ousi(n)} is the so-called movable
nu: it may appear or disappear according to phonetic and metrical context.

\subsection{The verb \emph{eimi}: to be}

Ancient Greek does not make the Portuguese lexical distinction between
\emph{ser} and \emph{estar} in this verb.  The present indicative of
\emph{eimi}, ``to be'', is irregular:
\begin{center}
\begin{tabular}{c|cc}
 & Singular & Plural \\
\hline
1st person & \emph{eimi} & \emph{esmen} \\
2nd person & \emph{ei} & \emph{este} \\
3rd person & \emph{esti(n)} & \emph{eisi(n)}
\end{tabular}
\end{center}
Hence Homeric \emph{eisi(n)} means ``they are''.  Its ending should not be
confused with the regular third-person plural active ending \emph{-ousi(n)}.

\subsection{A warning about apparent endings}

A sequence at the end of a word need not be a verbal ending.  For example,
Homeric \emph{alphestai}, found in the expression \emph{aneres alphestai}, is
a noun rather than a verb; it refers to men, especially men who earn or seek
their livelihood.  Its final \emph{-stai} should therefore not be parsed as a
verbal ending merely because it resembles one.

By contrast, verbal \emph{-tai} is a genuine third-person singular
middle/passive ending in several paradigms; an immediately preceding
\emph{s} may belong to the stem or tense formation rather than to the ending
itself.

\section{Ancient Greek numbers}

Some basic cardinal numbers are:
\begin{center}
\begin{tabular}{c|l}
1 & \emph{heis, mia, hen} \\
2 & \emph{duo} \\
3 & \emph{treis, tria} \\
4 & \emph{tettares/tessares} \\
5 & \emph{pente} \\
6 & \emph{hex} \\
7 & \emph{hepta} \\
8 & \emph{okto} \\
9 & \emph{ennea} \\
10 & \emph{deka} \\
100 & \emph{hekaton}
\end{tabular}
\end{center}

The word \emph{hecatomb} ultimately contains \emph{hekaton}, ``one hundred'',
and \emph{bous}, ``ox/cow'': historically a ``hundred-ox'' sacrifice, though
in Homer the term need not always imply a literal count of exactly one
hundred animals.

\section{Indo-European comparison}

The similarity of number words across Indo-European languages is often hidden
by regular sound change.  Greek \emph{hekaton} ``100'', Sanskrit \emph{satam},
Latin \emph{centum} (classically pronounced with initial /k/), and English
\emph{hundred} belong to the same inherited word-family, conventionally
reconstructed from Proto-Indo-European as approximately
\texttt{*kmtom}, with the initial consonant and syllabic nasal represented more
precisely in specialist notation.

The resemblance between Sanskrit \emph{eka}, ``one'', and part of Greek
\emph{hekaton}, ``one hundred'', is therefore accidental: Sanskrit
\emph{eka} is not the cognate of Greek \emph{hekaton}.  The Sanskrit cognate
of Greek \emph{hekaton} is \emph{satam}.  A more transparent comparison is
Greek \emph{deka}, ``ten'', with Sanskrit \emph{dasa}, ``ten''.

Polish belongs to the West Slavic branch of Indo-European:
Proto-Indo-European $\to$ Balto-Slavic $\to$ Slavic $\to$ West Slavic
$\to$ Polish.  Polish \emph{sto}, ``100'', belongs to the same inherited
family as the forms above.  English instead belongs to the Germanic branch:
Proto-Indo-European $\to$ Proto-Germanic $\to$ West Germanic $\to$
Anglo-Frisian $\to$ English.


\section{Etymological notes}

\subsection{\emph{Aiolos} and ``eolic''}

In the \emph{Odyssey}, Aiolos (Greek \emph{Aiolos}, Latin \emph{Aeolus}) is
the keeper of the winds: Zeus has given him power to restrain or release
them. This is his mythological office, however, not the literal etymological
meaning of his name. The name is related to the Greek adjective
\emph{aiolos}, whose range of meanings includes ideas such as quick-moving,
changeful, or shifting; its deeper etymology is not entirely certain.

The modern Romance adjective ``eolic'' (Portuguese \emph{eólico}, Spanish
\emph{eólico}, Italian \emph{eolico}, etc.) develops in the other direction:
Aiolos/Aeolus, because of his mythological association with the winds,
$\to$ ``pertaining to Aeolus'' $\to$ ``pertaining to the wind''. Thus
``eolic energy'' ultimately bears the name of Aiolos; Aiolos was not named
from a Greek word meaning ``wind''.

The Homeric association is especially vivid in \emph{Odyssey} 10.1--79. Odysseus and his companions reach the floating island of Aiolos, who entertains them for a month. When Odysseus asks to continue home, Aiolos gives him an ox-hide bag in which he has bound the courses of the troublesome winds, while he leaves the West Wind free to carry the ships toward Ithaca (10.19--27). After nine days they are already close enough to Ithaca to see the fires on shore, but Odysseus falls asleep. His companions, imagining that the bag contains gold and silver secretly given to their leader, open it; the imprisoned winds burst out and drive them all the way back to Aiolos (10.28--55). Aiolos then refuses to help Odysseus a second time, taking the disastrous return as evidence that the gods hate him (10.56--79). Thus the episode gives a concrete Homeric background for the later semantic development Aiolos/Aeolus $\to$ ``of Aeolus'' $\to$ ``of the wind''. 

The initial alpha in Greek \emph{Aiolos} carries a smooth breathing, not a
rough breathing. The smooth breathing has no consonantal phonetic
realization: there is no initial /h/. This contrasts with a rough breathing,
which does mark initial /h/. Thus the breathing sign is orthographically
present, but in this case it tells us precisely that the word begins directly
with the vowel.

\subsection{Homeric ``teeth-hedge'' and Icelandic \emph{tann-garðr}}

The following note is worth preserving verbatim, both for its Homeric
linguistic observation and for the comparison with Icelandic. It is Note 3
to p.~3, Book I.64, in S.~H. Butcher and A.~Lang, \emph{The Odyssey of Homer:
Done into English Prose} (Macmillan and Co., London, 1912; first printed
1879). The note itself credits E.~Magnússon, translator of \emph{Legends of Iceland}.

\begin{quote}
\emph{ἕρκος ὀδόντων}, or teeth-hedge.

This phrase, which regards the teeth as the hedge or fence which protects the
mouth, is a precise parallel to the Icelandic \emph{tann-garðr}, i.e.
teeth-garth (Old English \emph{geard}, enclosure). The difference, however,
between the two expressions is that the Icelandic is the common prose phrase
in use up to this day, while the Greek is archaic and poetical. The Icelandic
phrase is of too old a date in the language to be directly borrowed from the
Greek; it is the genuine metaphor of a military age, in which the teeth were
looked upon as the wall guarding the castle, that is, the mouth. A long list
might be drawn up of metaphorical expressions common to Homer and the
Icelandic sagas, but independent in origin, and pointing to similar customs
and conditions of life.

For this note we are indebted to Mr. E. Magnússon, translator of 'Legends of
Iceland,' etc.
\end{quote}

The Greek expression occurs at \emph{Odyssey} 1.64. The same note is also
cross-referenced by the editors at \emph{Odyssey} 10.328, in Circe's speech
after Odysseus drinks her \emph{pharmakon} without being enchanted: she says
that no other man has resisted the charm once it has ``passed his lips''.

\section{A first map of the Northern languages}

Here ``Northern languages'' will initially mean the North Germanic, or Nordic,
languages, while keeping their relation to the wider Germanic family visible:
\[
\text{Indo-European}\to\text{Germanic}\to\text{North Germanic}.
\]
The principal modern North Germanic languages are Icelandic, Faroese,
Norwegian, Danish, and Swedish. Old Norse is the medieval North Germanic
language of the sagas and much of the surviving Scandinavian literary
tradition. Icelandic is especially useful for reading this tradition because
its written language remains comparatively conservative, although modern
Icelandic is not simply identical with Old Norse.

This chapter will return later to the relations among Old Norse, Icelandic,
Old English, and the other Germanic languages, and to the literary question
that motivates the present note: why Northern languages and literatures were
so important to Borges. For now the Homer--Iceland comparison above is a
first concrete object to keep in view: similar metaphors may arise
independently, and resemblance by itself does not establish borrowing.
\bibliography{my}
\bibliographystyle{amsalpha}

\end{document}
